% Options for packages loaded elsewhere
\PassOptionsToPackage{unicode}{hyperref}
\PassOptionsToPackage{hyphens}{url}
\documentclass[
  ignorenonframetext,
]{beamer}
\newif\ifbibliography
\usepackage{pgfpages}
\setbeamertemplate{caption}[numbered]
\setbeamertemplate{caption label separator}{: }
\setbeamercolor{caption name}{fg=normal text.fg}
\beamertemplatenavigationsymbolsempty
% remove section numbering
\setbeamertemplate{part page}{
  \centering
  \begin{beamercolorbox}[sep=16pt,center]{part title}
    \usebeamerfont{part title}\insertpart\par
  \end{beamercolorbox}
}
\setbeamertemplate{section page}{
  \centering
  \begin{beamercolorbox}[sep=12pt,center]{section title}
    \usebeamerfont{section title}\insertsection\par
  \end{beamercolorbox}
}
\setbeamertemplate{subsection page}{
  \centering
  \begin{beamercolorbox}[sep=8pt,center]{subsection title}
    \usebeamerfont{subsection title}\insertsubsection\par
  \end{beamercolorbox}
}
% Prevent slide breaks in the middle of a paragraph
\widowpenalties 1 10000
\raggedbottom
\AtBeginPart{
  \frame{\partpage}
}
\AtBeginSection{
  \ifbibliography
  \else
    \frame{\sectionpage}
  \fi
}
\AtBeginSubsection{
  \frame{\subsectionpage}
}
\usepackage{iftex}
\ifPDFTeX
  \usepackage[T1]{fontenc}
  \usepackage[utf8]{inputenc}
  \usepackage{textcomp} % provide euro and other symbols
\else % if luatex or xetex
  \usepackage{unicode-math} % this also loads fontspec
  \defaultfontfeatures{Scale=MatchLowercase}
  \defaultfontfeatures[\rmfamily]{Ligatures=TeX,Scale=1}
\fi
\usepackage{lmodern}
\usecolortheme[]{dolphin}
\usefonttheme[]{structurebold}
\ifPDFTeX\else
  % xetex/luatex font selection
\fi
% Use upquote if available, for straight quotes in verbatim environments
\IfFileExists{upquote.sty}{\usepackage{upquote}}{}
\IfFileExists{microtype.sty}{% use microtype if available
  \usepackage[]{microtype}
  \UseMicrotypeSet[protrusion]{basicmath} % disable protrusion for tt fonts
}{}
\makeatletter
\@ifundefined{KOMAClassName}{% if non-KOMA class
  \IfFileExists{parskip.sty}{%
    \usepackage{parskip}
  }{% else
    \setlength{\parindent}{0pt}
    \setlength{\parskip}{6pt plus 2pt minus 1pt}}
}{% if KOMA class
  \KOMAoptions{parskip=half}}
\makeatother
\usepackage{graphicx}
\makeatletter
\newsavebox\pandoc@box
\newcommand*\pandocbounded[1]{% scales image to fit in text height/width
  \sbox\pandoc@box{#1}%
  \Gscale@div\@tempa{\textheight}{\dimexpr\ht\pandoc@box+\dp\pandoc@box\relax}%
  \Gscale@div\@tempb{\linewidth}{\wd\pandoc@box}%
  \ifdim\@tempb\p@<\@tempa\p@\let\@tempa\@tempb\fi% select the smaller of both
  \ifdim\@tempa\p@<\p@\scalebox{\@tempa}{\usebox\pandoc@box}%
  \else\usebox{\pandoc@box}%
  \fi%
}
% Set default figure placement to htbp
\def\fps@figure{htbp}
\makeatother
% definitions for citeproc citations
\NewDocumentCommand\citeproctext{}{}
\NewDocumentCommand\citeproc{mm}{%
  \begingroup\def\citeproctext{#2}\cite{#1}\endgroup}
\makeatletter
 % allow citations to break across lines
 \let\@cite@ofmt\@firstofone
 % avoid brackets around text for \cite:
 \def\@biblabel#1{}
 \def\@cite#1#2{{#1\if@tempswa , #2\fi}}
\makeatother
\newlength{\cslhangindent}
\setlength{\cslhangindent}{1.5em}
\newlength{\csllabelwidth}
\setlength{\csllabelwidth}{3em}
\newenvironment{CSLReferences}[2] % #1 hanging-indent, #2 entry-spacing
 {\begin{list}{}{%
  \setlength{\itemindent}{0pt}
  \setlength{\leftmargin}{0pt}
  \setlength{\parsep}{0pt}
  % turn on hanging indent if param 1 is 1
  \ifodd #1
   \setlength{\leftmargin}{\cslhangindent}
   \setlength{\itemindent}{-1\cslhangindent}
  \fi
  % set entry spacing
  \setlength{\itemsep}{#2\baselineskip}}}
 {\end{list}}
\usepackage{calc}
\newcommand{\CSLBlock}[1]{\hfill\break\parbox[t]{\linewidth}{\strut\ignorespaces#1\strut}}
\newcommand{\CSLLeftMargin}[1]{\parbox[t]{\csllabelwidth}{\strut#1\strut}}
\newcommand{\CSLRightInline}[1]{\parbox[t]{\linewidth - \csllabelwidth}{\strut#1\strut}}
\newcommand{\CSLIndent}[1]{\hspace{\cslhangindent}#1}
\setlength{\emergencystretch}{3em} % prevent overfull lines
\providecommand{\tightlist}{%
  \setlength{\itemsep}{0pt}\setlength{\parskip}{0pt}}
\usepackage{bookmark}
\IfFileExists{xurl.sty}{\usepackage{xurl}}{} % add URL line breaks if available
\urlstyle{same}
\hypersetup{
  pdftitle={Bayesian Graphical Lasso for Estimating Brain Networks},
  pdfauthor={Liam Hanson, BIOS 731},
  hidelinks,
  pdfcreator={LaTeX via pandoc}}

\title{Bayesian Graphical Lasso for Estimating Brain Networks}
\author{Liam Hanson, BIOS 731}
\date{}

\begin{document}
\frame{\titlepage}

\begin{frame}{Overview}
\phantomsection\label{overview}
\begin{itemize}
\item
  Discuss why we might want to estimate a graph
\item
  Discuss the Bayesian Graphical Lasso Gibbs Sampler
\item
  Review results from a small simulation study
\end{itemize}
\end{frame}

\begin{frame}{Motivation}
\phantomsection\label{motivation}
We are interested in viewing the brain as a graph i.e.~a set of nodes,
say regions of the brain, and edges, the connections between these edges

\begin{figure}
\includegraphics[width=0.6\linewidth]{figures/13-Figure1.6-1} \caption{[@Hendricksen2015VisualizingDB]}\label{fig:unnamed-chunk-1}
\end{figure}

Different brains have different network structures and therefore
different local and global properties of these networks
\end{frame}

\begin{frame}{Motivation}
\phantomsection\label{motivation-1}
In general, healthy brains have a so-called ``small-world'' or
``scale-free'' structure.

Each node has few neighbors but it is easy for one group of nodes or
sub-network to talk to another (six degrees of separation).

This implies the existence of hub nodes, nodes through which other areas
of the network communicate.
\end{frame}

\begin{frame}{Motivation}
\phantomsection\label{motivation-2}
\begin{figure}
\includegraphics[width=0.6\linewidth,height=0.6\textheight]{figures/41583_2009_Article_BFnrn2575_Fig3_HTML} \caption{[@Bullmore2009]}\label{fig:unnamed-chunk-2}
\end{figure}

A contrasting example is in people with Schizophrenia. These individuals
tend to score lower on these small-world metrics.
\end{frame}

\begin{frame}{Graph Estimation}
\phantomsection\label{graph-estimation}
We don't observe these graphs directly

\begin{figure}
\includegraphics[width=0.6\linewidth,height=0.6\textheight]{figures/41593_2013_Article_BFnn3423_Fig1_HTML} \caption{[@Buckner2013]}\label{fig:unnamed-chunk-3}
\end{figure}

Instead, we have time series data for each voxel or region of interest
in the brain where the measured quantity is the so-called BOLD signal.

We will use the associations between these time series to estimate our
graph.
\end{frame}

\begin{frame}{Graph Estimation}
\phantomsection\label{graph-estimation-1}
A graph can be represented as an adjacency matrix: a symmetric matrix
\(A\) with \(a_{ij} = 1\) if an edge is present between node \(i,j\) and
\(a_{ij} = 0\) otherwise.

For our purposes, we assume that at time step (i.e.~scan) \(t\), the
value of the BOLD signal across all \(D\) nodes
\(Y_D \sim MVN(0, \Sigma)\).

Fact: The precision matrix \(\Omega = \Sigma^{-1}\) has the property
that its entries \(\omega_{ij}\) encode the partial correlation between
node \(i\) and node \(j\).

If \(\omega_{ij} = 0\), these nodes are conditionally independent.
Generally, we expect the true precision matrices to be sparse (recall
scale free organization).
\end{frame}

\begin{frame}{Graphical Lasso}
\phantomsection\label{graphical-lasso}
\begin{itemize}
\item
  To capture this sparsity, something like LASSO (Tibshirani 1996) could
  work.
\item
  We could fit successive pairwise LASSO's and regress the residuals to
  obtain partial correlations, but this isn't ideal
\item
  Instead, let's maximize the penalized Multivariate Gaussian
  log-likelihood with respect to the \(\Omega\):
\end{itemize}

\[\textrm{log (det }\Omega) - \textrm{tr}(\frac{S}{n}\Omega) - \rho||\Omega||_1\]
where \(S\) is the sample covariance matrix and \(||\Omega||_1\) is the
entry-wise \(l_1\)-norm of \(\Omega\).
\end{frame}

\begin{frame}{Graphical Lasso}
\phantomsection\label{graphical-lasso-1}
\begin{itemize}
\item
  This is the approach of Friedman, Hastie, and Tibshirani (2007)
\item
  At a high level, it recursively fits a modified LASSO in a blockwise
  fashion using sub-gradient methods and coordinate descent (gets around
  non-differentiability of \(l_1\) norm.)
\item
  The focus of my project was a Bayesian Gibbs-Sampler approach to the
  Graphical Lasso by Wang (2012) to estimate the aforementioned graphs
\end{itemize}
\end{frame}

\begin{frame}{Bayesian Graphical Lasso via Gibbs Sampling}
\phantomsection\label{bayesian-graphical-lasso-via-gibbs-sampling}
We have the likelihood and prior
\[p(\textbf{y}_t | \Omega) = MVN(\textbf{y}_t|0, \Omega^{-1} )\]
\[p(\Omega|\lambda) = C^{-1} \prod_{i<j}\textrm{DE}(\omega_{ij}|\lambda)\prod_{i=1}^p
\textrm{Exp}(\omega_{ii}|\frac{\lambda}{2})\cdot1_{\Omega\in M^+}\]

where \(M^+\) is the space of positive-definite matrices.

\begin{itemize}
\tightlist
\item
  DE is the double-exponential (Laplace) distribution used in classic
  Bayesian LASSO which induces potential shrinkage to 0.
\item
  Natural exponential prior on the diagonal elements i.e.~inverted
  variances.
\end{itemize}
\end{frame}

\begin{frame}{Bayesian Graphical Lasso}
\phantomsection\label{bayesian-graphical-lasso}
\begin{itemize}
\item
  Fact: The double exponential results from marginalizing over a mixture
  of normals with variances \(\tau_{ij}\) and exponential mixing
  density.
\item
  This results in the posterior distribution:
\end{itemize}

\[p(\Omega, \tau|\textbf{Y}, \lambda) \propto |\Omega|^{\frac{n}{2}}\textrm{ exp\{-tr }
(\frac{1}{2}S\Omega)\}\prod_{i<j}\{\tau_{ij}^{-\frac{1}{2}}\textrm{exp}(-\frac{\omega_{ij}^2}{2\tau_{ij}})
\textrm{exp}(-\frac{\lambda^2}{2}\tau_{ij})\}\]

\[\times\prod_{i=1}^p\textrm{exp}(\frac{\lambda}{2}\omega_{ij})\cdot1_{\Omega\in M^+}\]

\begin{itemize}
\tightlist
\item
  The restriction to the space of positive definite matrices makes this
  tricky, but we can work with it.
\end{itemize}
\end{frame}

\begin{frame}{Deriving Gibbs Updates}
\phantomsection\label{deriving-gibbs-updates}
For each update, partition \(\Omega\), \(S\), and
\(\mathcal{T} = \{\tau_{ij}\}_{\in D\times D}\)

\[\Omega = \begin{bmatrix}
\Omega_{11} & \boldsymbol{\omega}_{12} \\
\boldsymbol{\omega}_{12}' & \omega_{22} 
\end{bmatrix} \space S = \begin{bmatrix}
S_{11} & \boldsymbol{s}_{12} \\
\boldsymbol{s}_{12}' & s_{22} 
\end{bmatrix} \space \mathcal{T} = \begin{bmatrix}
\mathcal{T}_{11} & \boldsymbol{\tau}_{12} \\
\boldsymbol{\tau}_{12}' & 0 
\end{bmatrix} \space\]

where \(\Omega_{11}, S_{11}, \mathcal{T}_{11}\) are block matrices and
\(\boldsymbol{\omega}_{12}, \boldsymbol{s}_{12}, \boldsymbol{\tau}_{12}\)
are length \(D-1\) vectors containing the joint precisions, sample
covariances, and scaling factors for the current update.
\end{frame}

\begin{frame}{Deriving Gibbs Updates}
\phantomsection\label{deriving-gibbs-updates-1}
Then the conditional joint distribution of
\((\boldsymbol{\omega}_{12},\omega_{22})\) is
\[p(\boldsymbol{\omega}_{12},\omega_{22}|\Omega_{11},\mathcal{T},\textbf{Y},\lambda) \propto (\omega_{22} - \boldsymbol{\omega}_{12}' \Omega_{11}^{-1}\boldsymbol{\omega}_{12})^{\frac{n}{2}}\]

\[\times \textrm{exp}[-\frac{1}{2}\{\boldsymbol{\omega}_{12}'D_{\tau}^{-1}\boldsymbol{\omega}_{12} + 2s'_{12}\boldsymbol{\omega}_{12} + (s_{22}+\lambda)\omega_{22}\}]\]

by dropping out constant terms and using determinant properties.

We then make the change of variables
\[(\boldsymbol{\omega}_{12},\omega_{22}) \to (\boldsymbol{\beta} = \boldsymbol{\omega}_{12}, \gamma = \omega_{22} - \boldsymbol{\omega}_{12}' \Omega_{11}^{-1}\boldsymbol{\omega}_{12})\]
resulting in new joint conditional . . .
\end{frame}

\begin{frame}{Deriving Gibbs Updates}
\phantomsection\label{deriving-gibbs-updates-2}
\[(\gamma, \boldsymbol{\beta})|(\Omega_{11}, \mathcal{T}, \textbf{Y}, \lambda) \sim
\textrm{Gamma}(\frac{n}{2}+1, \frac{s_{22}+\lambda}{2}) \times 
\textrm{N}(-\textbf{C}s_{21}, \textbf{C})\]

where
\(\textbf{C} = \{(s_{22}+\lambda)\Omega_{11}^{-1} + D_{\tau}^{-1}\}^{-1}\)

To sample from the full conditional of \(\tau\), use:

Fact: Conditional posteriors of \(u_{ij} = \frac{1}{\tau_{ij}}\) are
independently inverse Gaussian with parameters
\(\mu'=\sqrt{\lambda^2 / \omega_{ij}^2}, \lambda'=\lambda^2\).

Finally, we sample from the full conditional of our tuning parameter
\(\lambda\) given a prior \(\lambda \sim \textrm{Gamma}(r, s)\):

\[\lambda \sim \textrm{Gamma}(r + D(D+1)/2, s + ||\Omega||_1/2)\]
\end{frame}

\begin{frame}{Final Algorithm}
\phantomsection\label{final-algorithm}
This leads to the algorithm for the Block Gibbs Sampler. Given current
\(\Omega\) and \(\tau\):

\begin{enumerate}
\tightlist
\item
  For \(i=1,...,D\)
\end{enumerate}

\begin{enumerate}
[(a)]
\tightlist
\item
  Partition \(\Omega,S\) and \(\mathcal{T}\) as shown
\item
  Sample
  \(\gamma \sim (\Omega_{11}, \mathcal{T}, \textbf{Y}, \lambda) \sim \textrm{Gamma}(\frac{n}{2}+1, \frac{s_{22}+\lambda}{2})\)
  and
  \(\boldsymbol{\beta} \sim \textrm{N}(-\textbf{C}s_{21}, \textbf{C})\)
\item
  Update
  \(\boldsymbol{\omega}_{12}=\beta, \omega_{22} = \gamma + \beta'\Omega_{11}^{-1}\beta\)
\end{enumerate}

\begin{enumerate}
\setcounter{enumi}{1}
\item
  Sample \(\lambda\) from
  \(\lambda \sim \textrm{Gamma}(r + D(D+1)/2, s + ||\Omega||_1/2)\)
\item
  For \(i<j\), sample \(u_{ij} \sim \textrm{Inv-Gau}(\mu',\lambda')\)
  and update \(\tau_{ij} = \frac{1}{u_{ij}}\)
\end{enumerate}
\end{frame}

\begin{frame}{Implementation and Simulation Study}
\phantomsection\label{implementation-and-simulation-study}
\begin{itemize}
\item
  We will compare the performance of the Bayesian Graphical Lasso to its
  frequentist counterpart with a focus on edge classification
\item
  For Graphical Lasso, we simply say an edge exists between node \(i\)
  and node \(j\) if \(\omega_{ij} \neq 0\). Any edges which shrunk to 0
  are said to not exist.
\item
  For the Bayesian Graphical Lasso, the probability of our posterior
  estimates being exactly 0 is 0, so we say an edge exists if a 95\%
  credible interval about the posterior mean of \(\omega_{ij}\) does not
  contain 0.
\item
  Finally, for each simulation, we train the Graphical Lasso with
  10-fold cross validation and use its estimates of \(\Omega\) and
  \(\lambda\) as our initial Gibbs values (cheating).
\end{itemize}
\end{frame}

\begin{frame}{Data}
\phantomsection\label{data}
\begin{itemize}
\item
  We have four simulation scenarios with \(D \in {25,50,100,150}\) and
  max timepoints \(T=200\).
\item
  To simulate functional MRI data, we generate scale-free networks with
  \(D\) nodes then randomly rewire each node them with a small
  probability to mimic real brain networks.
\item
  That graph is converted into an adjacency matrix and then each
  non-zero edge is randomly scaled to create a sparse precision matrix
  \(\Omega_{true} \in D \times D\)
\item
  We then sample \(T\) times from \(MVN(0,\Omega_{true})\) to generate
  our simulated functional data.
\end{itemize}
\end{frame}

\begin{frame}{Simulation Results}
\phantomsection\label{simulation-results}
\includegraphics[width=0.48\linewidth,height=0.48\textheight]{figures/graph_plot_D_25}
\includegraphics[width=0.48\linewidth,height=0.48\textheight]{figures/graph_plot_D_50}
\includegraphics[width=0.48\linewidth,height=0.48\textheight]{figures/graph_plot_D_100}
\includegraphics[width=0.48\linewidth,height=0.48\textheight]{figures/graph_plot_D_150}
\end{frame}

\begin{frame}{Simulation Results}
\phantomsection\label{simulation-results-1}
Each simulation had 1000 samples with 200 burn-in samples. The sample
appeared to mix well quickly.

\pandocbounded{\includegraphics[keepaspectratio]{presentation_files/figures/diag_trace.png}}
\end{frame}

\begin{frame}{Simulation Results: Sampling}
\phantomsection\label{simulation-results-sampling}
\pandocbounded{\includegraphics[keepaspectratio]{presentation_files/figures/off_diag_trace.png}}
\end{frame}

\begin{frame}{Simulation Results: Sampling}
\phantomsection\label{simulation-results-sampling-1}
\pandocbounded{\includegraphics[keepaspectratio]{presentation_files/figures/zero_off_diag_trace.png}}
\end{frame}

\begin{frame}{Simulation Results: Edge Classification}
\phantomsection\label{simulation-results-edge-classification}
\pandocbounded{\includegraphics[keepaspectratio]{presentation_files/figures/sens.png}}
\end{frame}

\begin{frame}{Simulation Results: Edge Classification}
\phantomsection\label{simulation-results-edge-classification-1}
\pandocbounded{\includegraphics[keepaspectratio]{presentation_files/figures/spec.png}}
\end{frame}

\begin{frame}{Simulation Results: Stein Loss}
\phantomsection\label{simulation-results-stein-loss}
\pandocbounded{\includegraphics[keepaspectratio]{presentation_files/figures/stein.png}}
\end{frame}

\begin{frame}{Simulation Results: Coverage}
\phantomsection\label{simulation-results-coverage}
\pandocbounded{\includegraphics[keepaspectratio]{presentation_files/figures/omega_cov.png}}
\end{frame}

\begin{frame}{Simulation Results: Lambda}
\phantomsection\label{simulation-results-lambda}
\pandocbounded{\includegraphics[keepaspectratio]{presentation_files/figures/lambda.png}}
\end{frame}

\begin{frame}{Simulation Results: Computation Time}
\phantomsection\label{simulation-results-computation-time}
\pandocbounded{\includegraphics[keepaspectratio]{presentation_files/figures/comp.png}}
\end{frame}

\begin{frame}{Challenges}
\phantomsection\label{challenges}
\begin{itemize}
\item
  Beastly computation time for D \textgreater{} 150. Brain data even at
  coarse parcellations has higher dimension.
\item
  Extremely small values for \(\tau\) matrices challenges invertibility
\end{itemize}
\end{frame}

\begin{frame}{Future Directions}
\phantomsection\label{future-directions}
\begin{itemize}
\item
  Speed up computation dramatically, probably via Rcpp, so that this can
  be applied to real imaging data with higher dimension.
\item
  Implement Bayesian Classification rule
\item
  Bayesian Adaptive Graphical Lasso: allows for different \(\lambda\)'s
  at different edges.
\end{itemize}
\end{frame}

\begin{frame}{References}
\phantomsection\label{references}
\phantomsection\label{refs}
\begin{CSLReferences}{1}{0}
\bibitem[\citeproctext]{ref-friedman2007sparseinversecovarianceestimation}
Friedman, Jerome, Trevor Hastie, and Robert Tibshirani. 2007. {``Sparse
Inverse Covariance Estimation with the Lasso.''}
\url{https://arxiv.org/abs/0708.3517}.

\bibitem[\citeproctext]{ref-10.1214ux2f12-BA729}
Wang, Hao. 2012. {``{Bayesian Graphical Lasso Models and Efficient
Posterior Computation}.''} \emph{Bayesian Analysis} 7 (4): 867--86.
\url{https://doi.org/10.1214/12-BA729}.

\end{CSLReferences}
\end{frame}

\end{document}
